\section{The coefficient field at the last grouping scale}
\label{sec:ope-local-projection}
We connect the local-field representation to the degree-bounded
coefficient field in the cutoff estimate. This step uses a projection
estimate for one elementary insertion. We state its content explicitly
and prove the extension to the whole grouped observable, including
the normalization that determines its final price.

Let $m+1$ and $m$ be adjacent fine and coarse scales, with mesh
sizes $a_{m+1}=a_m/L$ and $L\ge1$. Fix a degree $d\ge0$.
For every fine elementary insertion $M$ anchored at the origin,
assume there is a coarse coefficient field $Q_M$ of degree at most
$d$ with the prescribed locality and covariance format, satisfying
\begin{equation}\label{eq:ope-leaf-projection}
 \operatorname{dens}(Q_M)=L^d\mathcal T\operatorname{dens}(M),
 \qquad |Q_M|_{\rm coeff}\le x\operatorname{cost}(M),
 \qquad \operatorname{rem}(Q_M)\le x\operatorname{cost}(M),
\end{equation}
where $x\ge0$. Here $\operatorname{rem}(Q)\le B$ means that $Q$ has
a regular remainder with summable majorants of total mass $R$ and a forced
remainder with uniformly summable majorants of total mass $F$ such that
$R+\kappa_m^{-1}F\le B$, where $\kappa_m>0$ is the forced-remainder
allowance of this scale; it is the normalization of the forced part used
in Section~\ref{sec:leaf-transfer}. The same $\kappa_m$ is used for
every field at this scale, in particular in the finite sums below. Assume, in addition, the singleton transition
expansion of Section~\ref{sec:local-transfer}: the fine expansion
cost for this insertion is at most $2\operatorname{cost}(M)$,
and its coarse response cost is at most
$x\operatorname{cost}(M)$. The relation for the quoted profile
has multiplier $L^d$, consistently with the first identity in
\eqref{eq:ope-leaf-projection}. These statements concern the
actual background transition and its actual coefficient fields;
they are not an abstract polynomial approximation assumption.
Section~\ref{sec:leaf-transfer} derives both statements from a localized
expansion of one real insertion under the block integral: any $Q_e$
satisfying \eqref{eq:leaf-projection-allowance} serves as $x$, with
$\operatorname{cost}(M)$ the fine mass $m$ there, fine expansion cost
$2m$ and coarse response cost $Q_em$. That derivation is conditional on
the localized expansion and the tail, flag and smallness inputs stated
in Section~\ref{sec:leaf-transfer}.
Section~\ref{sec:ope-history-service} places the projection on the
selected history with allowance $X_m=\exp(\sqrt{\widehat\beta_{m+1}})$,
for $m\ge t_d$, where $t_d$ is the reference level with $d+3\le J(t_d)$,
for every cutoff $N_0>m$ and every period.
The analytic inputs of Section~\ref{sec:leaf-transfer} remain open
here.

\subsection{From one insertion to a grouped field}
Suppose a grouped fine field has zero separately recorded constant
and an elementary representation
\[
        (g,q)=\sum_j z_j\bigl(\operatorname{ord}(M_j),
                              \operatorname{flag}(M_j)\bigr),
       \qquad \sum_j|z_j|\operatorname{cost}(M_j)\le A,
\]
where the list is finite and every insertion is anchored at the
origin. Choose $Q_j$ by \eqref{eq:ope-leaf-projection}, and let
$Q=\sum_jz_jQ_j$. The sum means addition of coefficient vectors
and of forced arrays, and concatenation of the regular remainder
lists. Thus all locality and covariance conditions are preserved.
The triangle inequality gives
\begin{equation}\label{eq:ope-group-projection}
       |Q|_{\rm coeff}\le xA,\qquad
       \operatorname{rem}(Q)\le xA.
\end{equation}
For the regular remainder, multiply each majorant by $|z_j|$
and sum its weighted costs. For the forced remainder, do the
same for its uniformly summable majorants. These constructions
prove the asserted bound without any cancellation of signed
terms. Finite linearity of integration gives
\[
          \operatorname{dens}(Q)
             =L^d\mathcal T\operatorname{coin}_R(g,q).
\]
Indeed the linear density of $(g,q)$ is the sum of the elementary
densities with coefficients $z_j$, by the singleton formula
\eqref{eq:ope-single-coin}. Each elementary density is required to be integrable under the
fine Haar prior and in each transition slice, as in
Section~\ref{sec:local-transfer}; the lawful representation and
transition-kernel properties supply this in the application. Thus
finite summation commutes with the transition integral. The empty list gives $Q=0$ and the same
conclusion.

The singleton transition expansions concatenate in precisely
the same way. Their fine cost is at most $2A$ and their coarse
response cost at most $xA$. Dividing the output field and its
profile by $L^d$ changes the latter cost to $L^{-d}xA$, while
the fine cost stays $2A$. The multiplier in the quoted-profile
identity is now one. This is the normalization needed by the
one-scale transport lemma.

Require also that the whole marked input density is integrable
under the fine prior and the transition slices; this is the
hypothesis needed below for taking polynomial coefficients and
regrouping through integration. Suppose the other grouped fields
are translated, coordinate
transformed or conjugated copies of $(g,q)$, with origins that
are embeddings of their coarse origins. The transition commutes
with these motions, and the coefficient-field format is stable
under them. Apply the same motion to $L^{-d}Q$ in each cell.
The quoted-profile expansions transform by relabelling the finite
set of slots; their cost bounds remain $2A$ and $L^{-d}xA$.
The singleton coefficient of the integrated marked density is
the linear density of the corresponding transformed copy of
$L^{-d}Q$: taking a coefficient commutes with integration, and
the preceding density identity commutes with the motion. Thus
the singleton-identification hypothesis of
Section~\ref{sec:local-transfer} holds for every copied cell,
with one and the same coefficient field $Q$.

If the incoming mark weights $I_v$ satisfy
\begin{equation}\label{eq:ope-projection-afford}
                 2I_v(2A)\le1,\qquad
                 I_vL^{-d}xA\le1,
\end{equation}
and the support, occupancy, numerical and new-weight hypotheses
of that transport lemma hold, it follows that the integrated
marked density has the required local representation with these
coefficient fields and higher-array price one. This conclusion
retains the entire marked density, not only its singleton
coefficients.

\subsection{The terminal price and physical normalization}
At the terminal grouping stage, let $z>0$ be any lower bound
for the common grouped weight supplied by the grouping induction.
Section~\ref{sec:ope-singletons} gives $A=4/z$. Take $x\ge1$ and
reduce the incoming weights to
\[
                             I_v=\frac{z}{16x}.
\]
They do not exceed the available grouped weight. Reducing positive
weights decreases every higher coefficient price. Moreover
\[
       2I_v(2A)=\frac1x\le1,
       \qquad I_vL^{-d}xA=\frac{L^{-d}}4\le1.
\]
This verifies \eqref{eq:ope-projection-afford}. If the outgoing
weights obey $0<J_v\le\delta_m I_v$, where $\delta_m$ is the
first reduction factor in the transport lemma, its final
reweighting condition follows term by term: the product of the
$J_v$ is at most the product of $\delta_m I_v$, and the allowed
support exponential is at least one.

We have obtained a coarse field $Q$ with both bounds at most
$4x/z$, appearing in the density with prefactor $L^{-d}$.
Define the physically normalized coefficient field
\[
                              Q'=a_{m+1}^dQ.
\]
Scaling multiplies its coefficient and remainder costs by the
positive number $a_{m+1}^d$. Since
\[
                     a_m^{-d}a_{m+1}^d=L^{-d},
\]
this changes neither the ordinary singleton nor its flagged
profile in the density. It only rewrites them with the standard
physical prefactor $a_m^{-d}$. We obtain
\begin{equation}\label{eq:ope-terminal-projection-price}
      |Q'|_{\rm coeff},\;\operatorname{rem}(Q')
                           \le a_{m+1}^d\frac{4x}{z}.
\end{equation}
This is the terminal projection price used in the cutoff
accounting: take $z$ equal to its available grouping budget and
$x$ equal to the elementary projection allowance.

Finally, regrouping and backwards integration are compatible.
The regrouping map reads the coefficient of the complete old
group at each new mark; it is a linear map between finite-dimensional
polynomial spaces. It therefore commutes with every integrable
transition. At the finest scale, reading a complete group in
the marked birth density gives the product of its constituent
observables. An induction over the intervening transitions shows
that the terminal reaped density is the backwards-integrated
density of that product. Consequently the coefficient field in
\eqref{eq:ope-terminal-projection-price} represents the intended
cluster product through the same sequence of scales. This
identifies the observable to which the cutoff price applies;
the elementary estimate \eqref{eq:ope-leaf-projection} and the
nonlinear transition estimate still require their own proofs.
